% Attempt at which each task was solved, per language (stacked bars, 10 tasks per language).
% Data: results/haiku-v1/report.md, section 4 (self-repair).
\begin{tikzpicture}
\begin{axis}[
  reportchart,
  xbar stacked,
  width=0.86\textwidth, height=6.2cm,
  bar width=10pt,
  xmin=0, xmax=10, xtick={0,2,...,10},
  xlabel={Zadania (z 10)},
  ytick={1,...,7}, ymin=0.4, ymax=7.6,
  yticklabels={Prolog, Scala, Haskell, Smalltalk, Java, Pascal, C},
  xmajorgrids=true,
  axis x line*=bottom, axis y line*=left,
  y axis line style={draw=none}, ytick style={draw=none},
  legend style={at={(0.0,1.02)}, anchor=south west, legend columns=4,
                /tikz/every even column/.append style={column sep=8pt}},
  every axis plot/.append style={draw=white, line width=0.8pt},
  legend image code/.code={\draw[#1, draw=none] (0cm,-0.1cm) rectangle (0.3cm,0.12cm);},
]
  \addplot[fill=seq600] coordinates {(10,7) (9,6) (6,5) (8,4) (7,3) (9,2) (7,1)};
  \addlegendentry{próba 0 (pierwsza odpowiedź)}
  \addplot[fill=seq400] coordinates {(0,7) (0,6) (2,5) (0,4) (2,3) (0,2) (1,1)};
  \addlegendentry{runda naprawy 1}
  \addplot[fill=seq250] coordinates {(0,7) (0,6) (2,5) (1,4) (0,3) (1,2) (0,1)};
  \addlegendentry{runda naprawy 2}
  \addplot[fill=neutral] coordinates {(0,7) (1,6) (0,5) (1,4) (1,3) (0,2) (2,1)};
  \addlegendentry{nierozwiązane}
\end{axis}
\end{tikzpicture}
